% !PRINTED 588
\weberpage{588}
%<<CONT
plexen Multiplication der elliptischen Functionen Beispiele giebt, ist das
Partial-Grundideal $\mathfrak{G}_R$ gleich 1 zu setzen.

Wenn die $n$ Körper (3) mit einander identisch sind, so heisst $\Omega$ ein
\gesperrt{Normalkörper in Bezug auf} $R$ (relativ normal, wenn $R$ nicht der absolute
Rationalitätsbereich ist). Die Gesammtheit der $n$ Substitutionen $(\Theta, \Theta_n)$,
die den Uebergang der Zahlen des Körpers $\Omega$ zu den conjugirten Zahlen vermitteln,
bilden eine Gruppe vom Grade $n$, wie wir schon im \S.\,147 des ersten Bandes gesehen
haben, die nichts anderes ist, als die \gesperrt{Galois'sche} Gruppe der Gleichung (1) im
Rationalitätsbereiche $R$, die wir hier die \gesperrt{Gruppe des Körpers $\Omega$ (in Bezug auf $R$)}
nennen.

Wenn insbesondere diese Gruppe commutativ ist, so ist $\Omega$ ein
\gesperrt{Abel'scher Körper in Bezug auf} $R$ (ein relativ \gesperrt{Abel'scher Körper}). In diesem
Falle ist (1) eine \gesperrt{Abel'sche} Gleichung im Rationalitätsbereiche $R$.

\begin{center}
\S.\,160.\\[6pt]
\gesperrt{Primideale im relativ normalen Körper.}
\end{center}

Wir nehmen jetzt an, dass der Körper $\Omega$ normal sei in Bezug auf den beliebigen
Körper $R$, und untersuchen unter dieser Voraussetzung die Primideale des Körpers
$\Omega$.

Die Gruppe von $\Omega$ in Bezug auf $R$ sei $\Phi$, und die Substitutionen von $\Phi$
mögen mit $\varphi, \varphi_1, \ldots$ bezeichnet sein. Wir bezeichnen ferner mit
\weq{1}{
\omega \mid \varphi
}
\gesperrt{die Zahl, die durch die Substitution $\varphi$ aus $\omega$ hervorgeht}, so dass $\omega$
und $\omega \mid \varphi$ in $\Omega$ enthalten sind.

Ebenso wie die Zahlen $\omega$ gehen auch alle Functionale und damit zugleich alle
Ideale $\mathfrak{a}$ des Körpers $\Omega$ durch eine Substitution $\varphi$ in bestimmte
Functionale und Ideale $\mathfrak{a} \mid \varphi$ über, und wenn $\omega$ eine durch
$\mathfrak{a}$ theilbare ganze Zahl ist, so ist $\omega \mid \varphi$ durch
$\mathfrak{a} \mid \varphi$ theilbar.

Wenn $\mathfrak{a}$ irgend ein Ideal in $\Omega$ ist, so giebt es gewisse Substitutionen
$\psi$ in $\Phi$, die der Bedingung
\weq{2}{
\mathfrak{a} \mid \psi = \mathfrak{a}
}
genügen, darunter immer die identische Substitution, und diese Substitutionen $\psi$
bilden eine Gruppe $\Psi$, die ein Theiler von $\Phi$
%CONT>>
